\section{The class of an elementary curve}\label{sec:local}

For a smooth genus-$g$ curve $C$ and $L\in\operatorname{Pic}^{g-1}(C)$
write $L_+=L$ and $L_-=K_C L^{-1}$. On a cycle of degree $m$ in
$\Tot(L_+\oplus L_-)$, evaluate the elementary symmetric functions of
the two fiber projections at $1+\max(0,m(g-1))$ distinct points of $C$.
The $j$th coefficients are holomorphic sections of $L_\pm^j$, including
at ramification by taking the characteristic polynomial on the finite direct
image from the normalization of each irreducible component and raising
it to the multiplicity of that component. Their vanishing at the chosen points implies
their identical vanishing for $1\le j\le m$. Thus the common zero set is
$m[C]$. In chosen frames the evaluations form an equivariant tuple $F$
of nonzero weights $r_1,\ldots,r_N$. On the associated bundle over
$\PP^N$ it defines a section of the pullback of $\bigoplus_j\cO(r_j)$.
The class
\[
 u_N=(\textstyle\prod_jr_j)^{-1}F^*\operatorname{Th}
                     \bigl(\textstyle\bigoplus_j\cO(r_j)\bigr)
\]
has compact support on the fixed section. Here $\operatorname{Th}$ is
the relative Thom class of the zero section. If $f$ is an additional test and
$i:\PP^N\hookrightarrow\PP^{N+1}$ is a hyperplane inclusion, then
$i_!u_N(F)=h\,u_{N+1}(F)=u_{N+1}(F,f)$, where $h=c_1(\cO(1))$.
The first equality is the hyperplane projection formula and the second
is multiplication by the normalized Euler class of the appended
line. Concatenation of two tuples therefore identifies their universal images
\cite[\S2]{PaperIII}. This construction defines $a_{g,m}$. It applies
also to continuous equivariant tests on an invariant neighborhood when
their common zero set is the same section.

\subsection{Deformation of the spin model}
Fix $L^2=K_C$ and a disk bundle $U$ about the zero section in
$L\oplus L$. Give the two summands the same Hermitian metric. On $U$
use a Kähler form invariant under real rotations of the copies of $L$.
In a spin frame $e^2=dz$ the holomorphic two-form in the construction of
\cite[Proposition~3.2]{IP} is
\[
                  \alpha=dz\wedge(w_1\,dw_1+w_2\,dw_2).
\]
It is the sum of the pullbacks of the tautological two-forms on the
two line summands. The metric and $\alpha$ are rotation invariant.
Hence the Lee almost-complex structures $J_{t\alpha}$ are invariant
for $0\le t\le T$ with $T>0$ small. On the two eigenlines
$w_1\pm iw_2$ the rotations have weights $+1,-1$.

\begin{proposition}\label{prop:local}
For $m>0$, the image of $a_{g,m}$ in $S$ equals the characteristic
function of the single cycle $m[C]$ in the elementary germ.
\end{proposition}
\begin{proof}
Let $Z_m$ be the integrable degree-$m$ cycle space in $U$, and let
$W_m$ be the corresponding space over the Lee interval. Each nonconstant
compact component has positive degree over $C$ by tameness and the
retraction of $U$ onto $C$. Image localization
\cite[Proposition~3.2]{IP} places every such component on $C$ for
$t>0$. Thus $W_m$ is the union of $\{0\}\times Z_m$ and
$[0,T]\times\{m[C]\}$, glued at $(0,m[C])$. Both subsets are closed
in the topology of parameter, support and current. Collapsing the
interval gives a continuous proper equivariant retraction onto $Z_m$.
Its inverse image of a compact set is that set together with the
interval if the set contains $m[C]$.

Pass to the associated bundles for the Hopf circle bundle over $\PP^N$.
At $t=0$ the unit-sphere reduction identifies this with the holomorphic
Borel family restricted to the disk. The pullback of $u_N$ under the
proper retraction has compact support on $\PP^N\times[0,T]$. At $t>0$
it is the normalized Euler class
$(\prod_jr_j)^{-1}\prod_j(r_jh)=h^N$, where $h=c_1(\cO(1))$.
Its integral is $+1$. By the interval relation of \S\ref{sec:families}
the two endpoints give the same class. At the positive endpoint the
Gysin map from a point of $\PP^N$ identifies $h^N$ with the unit at
$m[C]$ in the corresponding target fiber. Excision at the fixed section
permits any smaller disk. The retraction commutes with hyperplane
inclusion, so stabilization respects the entire interval.
\end{proof}

The class in $H$ itself is independent of the balanced splitting.
In a relative Picard chart choose the marked evaluation points as
disjoint local sections and trivialize the evaluation lines. The
fiberwise degree argument gives the same compact section of zeros
over a closed path segment. Its relative Thom class identifies the
endpoints in the analytic diagram under the interval relation.
The variety $\operatorname{Pic}^{g-1}(C)$ is connected. To vary $C$,
take successive charts in connected Teichmüller space, with the
relative degree-$(g-1)$ bundle $\cO((g-1)p)$ for a local marked section
$p$. The fixed-curve comparison identifies choices at common endpoints.
Use piecewise real-analytic paths in the charts, so each segment has
an analytic thickening. For genus zero use the constant sphere and for
genus one use marked elliptic curves. This also joins different spin
choices through balanced splittings.

\subsection{General elementary neighborhoods}
In an elementary neighborhood in the sense of
\cite[Definition~3.1]{IP}, the core is the only irreducible curve,
and its normal operator is the double of a real Cauchy--Riemann
operator $D_0$ on a complex line. Every branched-cover pullback of
this operator is injective. These conditions suffice to identify the
point insertion of the given germ with the class in
Proposition~\ref{prop:local}.

Let $C$ be the core. Choose a tubular chart of its normal bundle
$N$ adapted to the complex splitting along $C$, and pull back
the almost-complex structure $J$ to that chart. On a smaller fixed
disk use $J_s=\delta_s^*J$, $0<s\leq1$, where $\delta_s$ scales
the normal fibers by the positive real number $s$. The family
extends smoothly to the homogeneous almost-complex structure
$J_{\mathfrak d}$ determined by the normal Cauchy--Riemann
operator $\mathfrak d$. Indeed the only block in the pullback with
a negative power of $s$ is $s^{-1}B(x,sv)$ from tangent to normal. Since
$B(x,0)=0$, its extension at zero is obtained by integrating
$\partial_v B(x,\tau sv)[v]$ over $0\leq\tau\leq1$. Identification
of the limit with $J_{\mathfrak d}$ is the normal rescaling formula.
For $s>0$ the elementary isolation condition gives only the cycles
$m[C]$ in this disk. At zero an irreducible curve off the core
would project to a branched cover $f$ of $C$ with a nonzero section
in $\ker(f^*\mathfrak d)$, contrary to super-rigidity. Thus for
each $m>0$ the singleton insertion continues across this ordinary
interval~\cite[Lemmas~3.6 and~3.10 and Proposition~3.12]{DW} and
\cite[Definition 3.1]{IP}.

By the elementary splitting, $\mathfrak d=D_0\oplus D_0$ on
$L\oplus L$, with $D_0$ a real Cauchy--Riemann operator on a
complex line of degree $g-1$. Identify its underlying smooth
complex line bundle with that of a holomorphic spin line. On
this bundle join $D_0$ to the spin Dolbeault operator by affine
interpolation of the zero-order terms. The corresponding normal
models $J_{D_t\oplus D_t}$ project holomorphically to $C$, induce
the usual structure on each fiber and commute with real fiber
dilations and real rotations of the copies. These properties follow
from the construction of $J_{\mathfrak d}$ from the operator. A lift over a holomorphic $f$ is a pair of sections of
$f^*L$ satisfying $f^*D_t$. All models agree as almost-complex
structures along $C$. Thus one fixed Kähler form $\omega$ for the spin model
tames the whole path on a sufficiently small closed disk, by
compactness of $C\times[0,1]$.

Fix $m>0$ and let $r(Z)$ be the maximum fiber norm on the support
of a cycle, using a rotation-invariant metric. For a radius $a>0$
in the interior of the disk the locus $T_m$ of all $r(Z)=a$,
including the parameter $t$, is compact. Its supports lie in a
fixed compact set, and the taming form has integral precisely
$m\int_C\omega$ on each such cycle. Geometric compactness and
continuity of the radius give the claim. There are no circle-fixed
cycles in $T_m$. Each nonconstant irreducible component projects
to a finite holomorphic cover, since a fiber contains no compact
curve. An invariant cycle would therefore have a finite invariant
set in every fiber. The connected rotation group fixes every
element of such a set, and the origin is the only fixed vector~\cite[Theorem 2.11]{DIW}.

On a compact Hausdorff circle space without fixed points there
are finitely many continuous functions of nonzero weights having
no common zero. At a point $x$ its stabilizer is $\mu_k$ for some
$k\geq1$. The function $gx\mapsto g^k$ on its closed orbit
extends continuously to the space by Tietze extension applied to
the real and imaginary parts. Its weight-$k$ Fourier projection
agrees with the original function on the orbit. The nonvanishing
loci so obtained cover the space, hence a finite set suffices.
Apply this to $T_m$. For a test $f$ there, define a test on the
whole degree-$m$ cycle family by

\[
 \widehat f(Z)=
 \begin{cases}
 r(Z)\,f\bigl(t,(\delta_{a/r(Z)})_*Z\bigr),&r(Z)>0,\\
 0,&r(Z)=0 .
 \end{cases}
\]

Homogeneity ensures that the rescaled cycle lies in $T_m$.
Away from zero the formula is continuous by geometric convergence.
At zero it is continuous because $f$ is bounded. The tests retain
their weights and vanish simultaneously precisely at $m[C]$ in
each fiber. If $T_m$ is empty the family already consists of the
singleton in each fiber, and the empty tuple suffices.

Take Borel stage $N$ equal to the number of tests and pull back
their normalized relative Thom class. It has compact support on
the distinguished section over $\mathbb P^N\times[0,1]$. At the
super-rigid end its restriction is $h^N$ and gives the unit at
$m[C]$. At the integrable spin end it gives $a_{g,m}$ because the tests
are continuous circle-equivariant functions, and concatenation
with the holomorphic evaluation tuple after hyperplane stabilization
identifies their normalized Thom classes by the hyperplane formula.

The ordinary rescaling and this Borel deformation give the asserted
identity for the original germ. Write $b_{g,m}=\Phi(a_{g,m})$ and
$b_{g,0}=1$. On the point representatives, pullback by addition counts
the ordered pairs $(r[C],s[C])$ with $r+s=m$. Pullback by multiplication
by $d$ selects $(m/d)[C]$ when $d\mid m$ and is zero otherwise.
Both pullbacks preserve compact supports on these families. For addition,
the support and mass of either effective summand are bounded by those
of the sum, so compactness follows from \cite[Theorem~2.11]{DIW}.
The same reasoning applies to multiplication by $d$.
On these representatives the diagonal of the point base is a point.
The summands are external products of degree-zero classes, giving
\[
  \Delta b_{g,m}=\sum_{r+s=m}b_{g,r}\otimes b_{g,s},
  \qquad
  \rho_d b_{g,m}=
       \begin{cases} b_{g,m/d}&d\mid m,\\ 0&d\nmid m .\end{cases}
\]
